\documentclass[a4paper,11pt,french]{article}
\usepackage[utf8]{inputenc}

\usepackage[T1]{fontenc}
\usepackage[francais]{babel} 
\usepackage[top=2cm, bottom=2cm, left=2cm, right=2cm, includeheadfoot]{geometry} %pour les marges
\usepackage{lmodern}
\usepackage{pict2e}
\usepackage{fancyhdr} % Required for custom headers
\usepackage{lastpage} % Required to determine the last page for the footer
\usepackage{extramarks} % Required for headers and footers
\usepackage{graphicx} % Required to insert images
\usepackage{tabularx, longtable}
\usepackage{color, colortbl}
\usepackage{lscape}
%\usepackage[hidelinks]{hyperref}
\usepackage{longtable}
\usepackage{multirow}
\usepackage{rotating}
%\usepackage{pgfgantt}
%\usepackage{pgfcalendar}
%\usepackage{ifthen}
\usepackage{gensymb}


\linespread{1.1} % Line spacing

% Set up the header and footer
\pagestyle{fancy}
\lhead{\textbf{\hmwkClass -- \hmwkSubject \\ \hmwkTitle \\ \hmwkDocName}} % Top left header
\rhead{\includegraphics[width=10em]{../../images/logo_univ.png}}
\lfoot{\lastxmark} % Bottom left footer
\cfoot{} % Bottom center footer
\rfoot{Page\ \thepage\ / \pageref{LastPage}} % Bottom right footer
\renewcommand\headrulewidth{0.4pt} % Size of the header rule
\renewcommand\footrulewidth{0.4pt} % Size of the footer rule

\setlength{\headheight}{40pt}

\newcommand{\hmwkTitle}{Transchiffrement} % Assignment title
\newcommand{\hmwkClass}{Master 2 SSI } % Course/class
\newcommand{\hmwkAuthorName}{Émile GÉNÉRAT} % Your name
\newcommand{\hmwkSubject}{Conduite de projet} % Subject
\newcommand{\hmwkDocName}{Compte rendu réunion 30/01/14} % Document name

\makeatletter
\newcommand{\resettranslate}{\let\translate\@firstofone}
\makeatother

\definecolor{gris}{rgb}{0.95, 0.95, 0.95}

\title{
\vspace{2in}
\textmd{\textbf{\hmwkClass :\ \hmwkTitle}}\\
\normalsize\vspace{0.1in}\small{Due\ on\ \hmwkDueDate}\\
\vspace{0.1in}\large{\textit{\hmwkClassInstructor\ \hmwkClassTime}}
\vspace{3in}
}

\author{\hmwkAuthorName}
\date{} % Insert date here if you want it to appear below your name

\usepackage{amsmath}
\begin{document}
\newcount\startdate
\newcount\daynum
%\pgfcalendardatetojulian{2013-01-021}{\startdate}
\pagestyle{fancy}

\vspace*{5cm}
\begin{center}\textbf{\Huge{\hmwkDocName}}\end{center}
\vspace*{4.5cm}

\newpage

%Tableau de mises à jour


%La table des matières
%\clearpage
%\tableofcontents
%\clearpage

\section{Personnes présentes :}
\begin{itemize}
  \item Magali Bardet
  \item Emile Générat
  \item Julien Bourdon
  \item Yves Nouafo
  \item Ouissem Hamdani
  \item Jean-Baptiste Souchal
\end{itemize}

\section{Sujets abordés :}
Nous avons des problèmes pour établir une connexion TLS entre le client et le proxy.
Pour Madame Bardet, les points à vérifier/suggestions sont :
\begin{itemize}
\item s'assurer que le keystore contient les bonnes clés/certificats;
  \item le fichier cert.pem, qui doit contenir la concaténation des autorités de la chaîne;
  \item utiliser deux keystores;
  \item utiliser l'autorité intermédiaire faite en TP (pour ne pas avoir la clé privée de l'autorité racine) ;
  \item le certificat généré par le proxy pour un site doit avoir le même CN que le certificat du site distant.
  \item utiliser s\_client (si possible de le configurer avec un proxy)  
\end{itemize}

Pour l'acceptation des certificats d'autorité au sein des navigateurs, il faut chercher une solution sans devoir cliquer (problématique d'une flotte importante).

Nous avons pris du retard sur la partie Proxy. Madame Bardet nous demande de privilégier cette partie, qui est essentielle au projet.

\section{Actions à prévoir pour la réunion de la semaine prochaine :}
\begin{itemize}
  \item Présenter notre projet à Madame Bardet mardi matin, pour vérifier si la partie cryptographique est correcte.
\end{itemize}
~\\
La prochaine réunion est prévue le jeudi 13 février dans l'après-midi.

\end{document}